\section{Related Work}\label{sec:relwork}

  We now review techniques aimed at optimizing PCGs. Excluding recent works based on Ring LPN~\cite{ringLpn,foleage,pcgAnyField}, most PCGs rely on essentially the same protocols~\cite{CCS:BCGI18,silentOTCCS,silentOTcrypto,ringLpn,silver,EA} \\\noindent \cite{EC,structuredlpn}: they generate a secret sharing of the product \(\share{\Delta \mathbf{a}'}\), where \(\mathbf{a}'\in \Fbb^n\) is a sparse vector and \(\Delta\in \Fbb\) is a scalar. This corresponds to the interactive setup \(\mathbf{c}' - \mathbf{b}' = \Delta \mathbf{a}'\) (see \sectionref{sec:intro}). The resulting shared vector is then compressed using a linear function \(\mathbf{G}\). Consequently, improving PCG performance centers on two tasks: (1) generating the shared sparse vector \(\share{\Delta \mathbf{a}'}\) efficiently, and speeding up the subsequent multiplication by \(\mathbf{G}\).

Two major optimizations have emerged for generating the secret sharing of \(\Delta \mathbf{a}'\). The first is the concept of a regular distribution for \(\mathbf{a}'\)~\cite{AFS05}, where $\av'$ is the concatenation of $t$ unit vectors of length $n/t$. That is, \(n\) is the length of \(\mathbf{a}'\) and  \(t\) is its Hamming weight. Regular noise replaced the more traditional Bernoulli and exact distributions~\cite{AFS05,C:HOSS18,CCS:DoeShe17,CCS:BCGI18,CCS:SGRR19,silentOTCCS}\\ \cite{silentOTcrypto,ringLpn,EA}. 
In the context of secure computation, this approach was first used by \(\mathsf{TinyKeys}\)~\cite{C:HOSS18} and subsequently by~\cite{CCS:BCGI18} for PCGs. 
The regular noise distribution is advantageous because it allows generating \(\share{\Delta \mathbf{a}'}\) with only \(O(n)\) AES calls, spread across \(t\) distributed point functions.
This is in contrast with the other noise distributions, which typically require \(O(t)\) distributed point functions, each of size $n$, resulting in \(O(tn)\) AES calls. Alternatively, one can use more complex techniques such as cuckoo hashing~\cite{CCS:SGRR19}, but these also impose additional overheads.

More recently, \cite{structuredlpn} presented a new assumption that allows the cost of generating \(\share{\Delta \mathbf{a}'}\) to be amortized across \(q\) instances by fixing the noisy locations of $\av'$. This allows for many correlated noise vectors $\av'^{(1)},...,\av'^{(q)}\in \Fbb^n$ to be generated more efficiently.
As long as the sparse noise elements in each $\av'^{(i)}$ are sampled uniformly at random, the resulting $\av^{(i)}:=\G\av'^{(i)}\in \Fbb^k$ will be pseudorandom by the stationary LPN assumption \cite{structuredlpn}.
This leads to a concrete \(6\textup{--}18\times\) reduction in communication along with compelling performance improvements for our GPU implementation.

Another line of research, which constitutes the principal runtime bottleneck for PCGs and is the focus of this work, aims to reduce the cost of multiplying \(\mathbf{G}\) by \(\share{\Delta \mathbf{a}'}\)~\cite{CCS:BCGI18,silentOTCCS,silentOTcrypto,ringLpn,silver,EA,EC}. 
In the original formulation of Syndrome Decoding and LPN, \(\mathbf{G}\) is taken to be uniform, so computing \(\mathbf{G} \cdot \share{\Delta \mathbf{a}'}\) requires \(O(n^2)\) time. 
However, there is good evidence that one may replace \(\mathbf{G}\) with the generator matrix of a code that has high minimum distance, unless that code exhibits a strong algebraic structure (e.g., a Reed--Solomon code with efficient Syndrome Decoding via Berlekamp--Massey~\cite{Berlekamp68}). %In particular, can be efficiently decoded by the Berlekamp--Massey algorithm~\cite{Berlekamp68}. 
Switching \(\mathbf{G}\) to a generator matrix of a linear-time code makes it theoretically possible to evaluate \(\mathbf{G} \cdot \mathbf{v}\) in \(O(n)\) time for any \(\mathbf{v}\). 
In practice, however, performance often degrades to \(O(c n)\), where \(c\) is typically related to \(\log n\) or the security parameter.
Several works~\cite{CCS:BCGI18,silentOTCCS,silentOTcrypto} propose using LDPC or quasi-cyclic codes, which offer well-studied security and reasonable performance. 
Subsequently, \cite{EA} introduced a turbo-like code, called Expand-Accumulate, which significantly reduces costs and can be parameterized to achieve a provably high minimum distance. 
Further improvements were proposed by~\cite{EC}, who replaced \cite{EA}'s accumulation step with a convolution, thereby boosting minimum distance and enabling more favorable parameters. 
However, none of these existing designs focus on codes suitable for efficient GPU acceleration. Moreover, existing proof techniques do not result in codes with optimal minimum distance. When translated into the context of PCGs, low minimum distance means that the setup phase must be more expensive to maintain security. 
Our primary objective, therefore, is to develop a code with extremely high minimum distance that can be implemented efficiently on \emph{both} CPUs and GPUs.

In the ZK literature, some additional works such as \cite{field,blaze} generalize RA codes \cite{spielman} to work over larger fields. We note that our codes could be generalized to the non-binary case using our core techniques. 